%% ============================================================
%% Explainable Decision Intelligence for Clinical Laboratory Operations
%% IEEE Access manuscript source (compiles on Overleaf with the
%% built-in IEEEtran class; no local package install required).
%%
%% To compile: upload this file to a new Overleaf project, set the
%% compiler to pdfLaTeX, and place the figure PDFs from
%% ../results/figures/ into the project.
%% ============================================================
\documentclass[journal]{IEEEtran}

\usepackage{graphicx}
\usepackage{amsmath,amssymb}
\usepackage{booktabs}
\usepackage{array}
\usepackage{url}
\usepackage{xcolor}
\usepackage[hidelinks]{hyperref}

\begin{document}

\title{Explainable Decision Intelligence for Clinical Laboratory Operations:
Predicting Turnaround-Time Breaches and Specimen Rejection for Cost-Aware
Quality Management}

\author{Jaisritha~Damavarapu% <-this % stops a space
\thanks{J. Damavarapu is with the Doctor of Business Administration program
(Emphasis: AI and Digital Transformation Strategy/Health Management Innovation), Midwest University, USA,
and is a practicing laboratory technologist.
E-mail: djaisritha@gmail.com.}%
\thanks{Reproducible code and data generator:
\url{https://github.com/jaisrithad/clinical-lab-decision-intelligence}.}}

\markboth{Preprint: submitted to Knowledge-Based Systems (Elsevier)}%
{Jaisritha Damavarapu: Explainable Decision Intelligence for Clinical Laboratory Operations}

\maketitle

\begin{abstract}
Clinical and diagnostic laboratories operate under intensifying pressure to
deliver fast, reliable results while controlling cost, yet the operational
decisions that govern turnaround time (TAT) and specimen quality are still
largely reactive and experience-driven. Prior machine-learning work predicts
laboratory TAT accurately but stops at prediction, framing the problem as a
clinical-efficiency task rather than a managerial decision problem. This study
develops and validates an explainable decision-intelligence framework that
couples predictive modeling with interpretability and an explicit cost-of-quality
optimization, so that model outputs become governed, economically rational
operating decisions. Using a reproducible, literature-calibrated dataset of
20{,}000 laboratory test orders, we benchmark logistic regression, random forest,
and gradient boosting on two leakage-safe tasks, TAT-breach prediction and
pre-analytical specimen rejection, under stratified cross-validation and a
held-out test set. The best model attains a test ROC-AUC of 0.886 for TAT-breach
prediction; rejection proves intrinsically harder (ROC-AUC 0.702), which is itself
an actionable finding. SHAP attributions identify workload-at-arrival, analyzer
utilization, request priority, night-shift operation, and transport time as the
dominant, managerially-actionable drivers. Translating predictions into
expedite-or-not decisions through a cost-sensitive threshold reduces expected
operational cost by 35.9\% relative to no intervention and by 23.7\% relative to
the best managerial heuristic, with the advantage widening to 50--78\% as the
stakes rise, as in time-critical donor-screening contexts. The framework offers
laboratory and diagnostic-service managers a transparent, auditable bridge from
prediction to decision, and contributes a replicable template for
decision-intelligence research in regulated biomedical operations.
\end{abstract}

\begin{IEEEkeywords}
Decision intelligence, explainable artificial intelligence, clinical laboratory
operations, turnaround time, cost of quality, healthcare analytics, SHAP,
digital transformation.
\end{IEEEkeywords}

\section{Introduction}
\IEEEPARstart{T}{he} clinical laboratory is the quiet backbone of modern
healthcare: an estimated majority of medical decisions depend on laboratory
results, and the speed and reliability with which those results are produced
shape patient flow, clinician scheduling, and the cost of care. Two operational
indicators dominate how laboratories are judged. The first is \emph{turnaround
time} (TAT), the interval from order or collection to a reported result, which
is the single most visible performance metric a laboratory reports to its
clinical customers. The second is \emph{specimen quality}, operationalized at the
front door of the laboratory as the pre-analytical rejection of samples that are
hemolyzed, clotted, mislabeled, or of insufficient volume. Both indicators carry
direct economic weight: a breached TAT can extend an emergency-department stay or
trigger duplicate testing, while a rejected specimen forces recollection, delays
diagnosis, and consumes staff time.

These are, at root, \emph{management} problems. Deciding when to expedite a
sample, how to staff a shift, and where to intervene in the workflow are
resource-allocation decisions made under uncertainty and cost constraints. Yet in
most laboratories these decisions remain reactive, triggered after a delay has
already occurred, and are rarely informed by predictive, auditable analytics.
This gap is strategically important in exactly the settings where the stakes are
highest. In time-critical donor screening for organ and tissue donation, for
example, serological and nucleic-acid testing sits on the critical path to organ
viability; a delayed or rejected specimen is not merely an efficiency loss but can
foreclose a transplant. The same analytics that optimize routine laboratory
throughput therefore carry disproportionate value where delay is catastrophic.

Machine learning has begun to enter this space, but narrowly. Existing studies
demonstrate that TAT can be predicted from operational features with useful
accuracy~\cite{tatml2022,pathoml2026}, and that data-mining pipelines can monitor
TAT and prompt targeted interventions~\cite{orange2026,urinepretat2025}. These
contributions are valuable, but they share two limitations. First, they are framed
as clinical-efficiency or laboratory-medicine problems and are published in that
literature, leaving the managerial and strategic dimension, how a prediction
becomes a cost-rational decision, largely unaddressed. Second, they stop at the
predictive step: a probability is produced, but the decision rule that converts it
into action, and the economic consequences of that rule, are not modeled.

This paper closes that gap by reframing laboratory operations as an
\emph{explainable decision-intelligence} problem. We make three contributions:

\begin{enumerate}
\item \textbf{A decision-intelligence framework, not merely a predictor.} We
couple predictive modeling with model-agnostic explainability (SHAP) and an
explicit cost-of-quality optimization, so that predictions are translated into
governed, economically rational operating decisions. This three-layer design is
the central methodological contribution and the point of departure from prior TAT
work.
\item \textbf{A rigorous, reproducible dual-task benchmark.} We formulate two
leakage-safe decision tasks, TAT-breach prediction and pre-analytical specimen
rejection, and benchmark three model families under stratified cross-validation
and a held-out test set, reporting discrimination, calibration, and
decision-relevant metrics. The literature-calibrated data generator and all code
are released, so every reported number is reproducible.
\item \textbf{An economic and managerial interpretation.} We quantify the cost
savings of model-guided intervention against managerial baselines, conduct a
sensitivity analysis over the cost structure, and connect the results to concrete
managerial levers and to high-stakes settings such as time-critical donor
screening.
\end{enumerate}

The remainder of the paper reviews related work (Section~\ref{sec:related}),
formulates the problem and describes the data (Section~\ref{sec:data}), details
the methodology (Section~\ref{sec:method}), presents predictive results
(Section~\ref{sec:results}) and the explainability and decision analysis
(Section~\ref{sec:decision}), discusses implications and limitations
(Sections~\ref{sec:discussion}--\ref{sec:limits}), and concludes
(Section~\ref{sec:conclusion}).

\section{Related Work}\label{sec:related}

\subsection{Machine learning for laboratory turnaround time}
TAT prediction has attracted growing attention as laboratories digitize. A study
on roughly ninety thousand clinical-chemistry samples from a European academic
medical center showed that ensemble regressors could predict TAT from features
such as request priority and the workload present at sample arrival, enabling
earlier action on at-risk samples~\cite{tatml2022}. Random-forest models have been
applied to anatomic-pathology TAT, a domain where processing is lengthy and
variable~\cite{pathoml2026}. Beyond point prediction, low-code data-mining
pipelines have been used to monitor TAT over time and to trigger process
interventions, with reported reductions in pre-analytical and analytical
intervals~\cite{orange2026}. Related work on pre-analytical TAT for microbiology
has combined regression with clustering to attribute delay to specific workflow
stages and ward contexts~\cite{urinepretat2025}.

These studies establish that laboratory operational outcomes are learnable from
routinely captured features. However, they are positioned within
laboratory-medicine and clinical-efficiency framing, optimize predictive accuracy
or monitoring, and do not model the downstream decision or its cost. The present
work builds directly on their feature intuitions, workload, priority, shift, and
specimen characteristics, while shifting the objective from prediction to
cost-rational decision-making.

\subsection{Pre-analytical quality and the total testing process}
The pre-analytical phase is the dominant source of error in the total testing
process, and specimen rejection is its most visible control
point~\cite{plebani2010,hawkins2012}. Standards bodies specify acceptance and
rejection criteria and collection-and-transport requirements~\cite{clsi2021}.
Rejection is comparatively rare and heterogeneous in its causes, which, as our
results confirm, makes it intrinsically harder to predict than TAT breach, an
asymmetry with practical consequences for where predictive effort is best spent.

\subsection{Cost of quality and decision support}
The cost-of-quality tradition in operations management distinguishes prevention
and appraisal costs from the internal and external failure costs they are meant to
avert~\cite{crosby1979}. Framing laboratory intervention decisions in these terms
converts a classification output into an economic optimization: the cost of
expediting (an appraisal/prevention action) is weighed against the expected
failure cost of an unmitigated breach. Decision-support and decision-intelligence
systems have a long history in management research, but their application to
laboratory operations with an explicit, explainable cost model is, to our
knowledge, undeveloped. We position our contribution at this intersection.

\subsection{Explainable AI in regulated settings}
In regulated, safety-relevant domains, a prediction that cannot be explained is
difficult to govern or to act upon. SHAP provides a theoretically grounded,
model-agnostic attribution of a prediction to its input features~\cite{shap2017},
and is well suited to translating a risk score into a managerial cause (for
example, attributing a flagged sample to excess workload-per-staff rather than to
specimen type). Explainability is therefore not cosmetic here; it is what makes the
decision layer auditable and actionable.

\section{Problem Formulation and Data}\label{sec:data}

\subsection{Decision tasks}
We study two binary decision problems defined at the moment a specimen is
accessioned, so that every feature is available at decision time and no label
leakage occurs:
\begin{itemize}
\item \textbf{Task A: Specimen rejection.} Will the specimen be rejected at
accessioning for a pre-analytical quality failure?
\item \textbf{Task B: TAT breach.} Among processed (non-rejected) specimens,
will total turnaround time exceed the clinically defined threshold for the
sample's priority class?
\end{itemize}
Priority-specific TAT thresholds are set at 60, 100, and 150 minutes for STAT,
urgent, and routine requests respectively, reflecting the common practice of
tighter targets for higher-acuity requests.

\subsection{Data and calibration}
Sample-level laboratory operational data, timestamps, workload, and specimen
quality, are not released publicly by laboratories or organ-procurement
organizations, for patient-privacy and commercial reasons, and credentialed
registries require data-use agreements that fall outside a short research cycle.
We therefore generate a synthetic dataset of 20{,}000 test orders whose marginal
distributions and effect sizes are calibrated to values reported in the
peer-reviewed literature, and we release the generator so that all results are
fully reproducible. The aim of calibration is not to reproduce any single real
laboratory but to ensure that the \emph{relationships} a model must learn resemble
those observed in practice.

Calibration targets are drawn from the literature reviewed in
Section~\ref{sec:related}: specimen rejection is concentrated in hemolysis-prone
specimen types and high-workload, night-shift, and emergency contexts, at an
overall rate (3.0\%) within the realistic band for a mixed-acuity laboratory; TAT
is right-skewed and driven primarily by the workload present at arrival, request
priority, shift, analyzer utilization, and add-on or repeat testing. The resulting
dataset exhibits a 3.0\% rejection rate and, among processed samples, a 5.2\%
TAT-breach rate (STAT 18.5\%, urgent 7.1\%, routine 1.2\%), yielding two realistic,
class-imbalanced decision problems. Each record carries operational-context
features (priority, laboratory section, specimen type, collection source, hour,
shift, weekend), load features (pending-sample queue, staff on duty,
workload-per-staff, analyzer utilization, add-on and repeat flags), and
pre-analytical quality signals (hemolysis index, transport time, label
completeness, volume adequacy).

\section{Methodology}\label{sec:method}

\subsection{Models and preprocessing}
We benchmark three model families spanning the interpretability--flexibility
spectrum: $\ell_2$-regularized logistic regression (LR), random forest (RF), and
gradient-boosted trees (XGBoost)~\cite{breiman2001,xgboost2016}. Categorical
features are one-hot encoded; numeric features are standardized for LR.
Class imbalance is addressed with balanced class weighting for LR and RF and with
a scale-positive-weight adjustment for XGBoost. All pipelines are implemented in
scikit-learn and XGBoost~\cite{sklearn2011}.

\subsection{Validation protocol}
For each task the data are split into a 75\% training set and a 25\% stratified
held-out test set. On the training set we estimate generalization with stratified
five-fold cross-validation, reporting the mean and standard deviation of ROC-AUC,
average precision (PR-AUC), and the Brier score. Final performance is reported on
the untouched test set. We emphasize three complementary metrics because they
answer different questions: ROC-AUC measures ranking discrimination, PR-AUC is
informative under class imbalance, and the Brier score measures probability
calibration, which matters directly for the cost-threshold decision that follows.

\subsection{Explainability}
For the best-performing TAT-breach model we compute SHAP values on the test
set~\cite{shap2017}, summarizing global importance as the mean absolute SHAP value
per feature. This attributes the model's risk scores to specific operational
drivers, each of which maps to a managerial lever.

\subsection{Cost-of-quality decision layer}
The decision is whether to \emph{intervene}, to flag and expedite a
sample, given its predicted breach probability $\hat{p}$ and a decision threshold
$t$: intervene if $\hat{p}\geq t$. Let $C_{\mathrm{I}}$ be the marginal cost of
intervening on one sample (staff time to flag and expedite) and $C_{\mathrm{B}}$
the expected downstream cost of one unmitigated breach (extended length of stay,
repeat orders, clinician follow-up). Assuming intervention mitigates a true
breach, the expected total cost at threshold $t$ over $N$ samples is
\begin{equation}
\mathcal{C}(t) = C_{\mathrm{I}}\bigl(\mathrm{TP}(t)+\mathrm{FP}(t)\bigr)
+ C_{\mathrm{B}}\,\mathrm{FN}(t),
\label{eq:cost}
\end{equation}
where $\mathrm{TP}$, $\mathrm{FP}$, and $\mathrm{FN}$ are the counts of true
positives, false positives, and false negatives induced by $t$. The cost-optimal
operating point is
\begin{equation}
t^{\star} = \operatorname*{argmin}_{t\in(0,1)} \mathcal{C}(t).
\label{eq:argmin}
\end{equation}
Because $t^{\star}$ is obtained by directly minimizing the empirical cost in
Equation~\eqref{eq:cost} over the observed probabilities, the procedure is robust
to imperfect probability calibration: it finds the best operating point for
whatever scores the model produces.

We compare the model-guided policy against three managerial baselines: (i)
\emph{treat none} (no intervention), (ii) \emph{treat all} (expedite every
sample), and (iii) a \emph{priority-only rule} (expedite STAT requests only), the
heuristic a laboratory would most plausibly use without analytics. Illustrative
unit costs ($C_{\mathrm{I}}=\$15$, $C_{\mathrm{B}}=\$120$) are used for the main
comparison, and a sensitivity analysis varies the cost ratio
$C_{\mathrm{B}}/C_{\mathrm{I}}$ from 2 to 50 to represent settings from routine
testing to time-critical donor screening.

\section{Predictive Results}\label{sec:results}
Table~\ref{tab:models} reports predictive performance for both tasks.
For TAT-breach prediction, all three models discriminate well; logistic regression
attains the highest test ROC-AUC (0.886) and PR-AUC (0.433), with cross-validation
tightly concentrated (0.894\,$\pm$\,0.008), indicating stable generalization.
Figure~\ref{fig:perf} shows the corresponding ROC and precision--recall curves.
The strong ROC-AUC alongside a far lower PR-AUC is the expected signature of an
imbalanced problem and underscores why average precision, not accuracy, is the
appropriate discrimination metric here.

A nuance with direct operational relevance appears in the calibration column:
although logistic regression ranks best by discrimination, the tree ensembles are
better \emph{calibrated} (Brier 0.070--0.072 versus 0.125). Where a laboratory
intends to act on raw probabilities rather than on a tuned threshold, the
better-calibrated model may be preferable; where, as in our decision layer, an
empirical cost-optimal threshold is selected, discrimination is the binding
consideration and the robustness noted after Equation~\eqref{eq:argmin} applies.

Specimen rejection is substantially harder to predict (best test ROC-AUC 0.702).
This is not a modeling deficiency but a property of the phenomenon: rejection is
rare and its causes are heterogeneous and partly stochastic. The practical
implication is managerial. Predictive effort and automation are better
concentrated on TAT management, while rejection is better addressed through
upstream process controls at collection and transport.

\begin{table*}[!t]
\centering
\caption{Predictive performance across models and tasks (mean $\pm$ s.d. over
5-fold cross-validation; test metrics on a held-out 25\% split). Brier is a
calibration metric (lower is better); F1 is reported at the 0.50 threshold.}
\label{tab:models}
\begin{tabular}{llccccc}
\toprule
Task & Model & CV ROC-AUC & Test ROC-AUC & PR-AUC & Brier & F1 \\
\midrule
Specimen rejection & Logistic Regression & 0.705 $\pm$ 0.031 & 0.702 & 0.172 & 0.187 & 0.139 \\
Specimen rejection & Random Forest & 0.714 $\pm$ 0.036 & 0.682 & 0.149 & 0.060 & 0.224 \\
Specimen rejection & XGBoost & 0.634 $\pm$ 0.035 & 0.654 & 0.093 & 0.069 & 0.144 \\
TAT breach & Logistic Regression & 0.894 $\pm$ 0.008 & 0.886 & 0.433 & 0.125 & 0.313 \\
TAT breach & Random Forest & 0.877 $\pm$ 0.010 & 0.874 & 0.385 & 0.072 & 0.390 \\
TAT breach & XGBoost & 0.864 $\pm$ 0.009 & 0.871 & 0.369 & 0.070 & 0.364 \\
\bottomrule
\end{tabular}
\end{table*}

\begin{figure*}[!t]
\centering
\includegraphics[width=0.92\textwidth]{figure1_performance.pdf}
\caption{Predictive performance on the held-out test set for the TAT-breach task.
(a) ROC curves and (b) precision--recall curves for the three models; the dashed
line in (b) marks the positive-class prevalence.}
\label{fig:perf}
\end{figure*}

\section{Explainability and Decision Analysis}\label{sec:decision}

\begin{figure*}[!t]
\centering
\includegraphics[width=0.95\textwidth]{figure2_decision_intelligence.pdf}
\caption{Decision-intelligence layer for the TAT-breach task. (a) Global SHAP
importance of the risk drivers (mean absolute SHAP value); the leading drivers
after request priority are operational variables a manager can influence. (b)
Expected total operational cost as a function of the decision threshold; the
convex curve attains its minimum at the cost-optimal threshold ($t^{\star}=0.62$,
marked), lying below both the treat-none and priority-only baselines.}
\label{fig:decision}
\end{figure*}

\subsection{What drives a breach}
Figure~\ref{fig:decision}(a) ranks the global SHAP importance of the TAT-breach
drivers. Request priority is dominant, unsurprising given priority-specific
thresholds, but the next tier is composed almost entirely of \emph{operational}
variables that management can influence: analyzer utilization, the pending-sample
queue at arrival, transport time, night-shift operation, add-on testing, and
workload-per-staff. This is the actionable core of the result: the model does not
merely flag risk, it attributes risk to levers a laboratory manager controls
through staffing, shift design, analyzer capacity planning, and transport
logistics.

\subsection{From prediction to cost-rational decision}
Figure~\ref{fig:decision}(b) plots expected total cost against the decision
threshold. The curve is convex with a clear minimum at $t^{\star}=0.62$, and over a
wide band of thresholds the model-guided policy lies below both the treat-none and
priority-only baselines. Table~\ref{tab:cost} quantifies the comparison on the
held-out test set. Doing nothing incurs the full failure cost; expediting
everything is worst of all, more than doubling cost; the priority-only heuristic
helps modestly (16.0\% below treat-none). The model-guided policy at its
cost-optimal threshold reduces expected cost by 35.9\% relative to treat-none and
by 23.7\% relative to the best managerial baseline, the economically meaningful
comparison, since no competent manager does nothing.

\begin{table*}[!t]
\centering
\caption{Cost-of-quality comparison of decision policies on the held-out test set
(unit costs: \$15 to intervene, \$120 per unmitigated breach). Savings are
relative to doing nothing and to the best non-model baseline.}
\label{tab:cost}
\begin{tabular}{lccccc}
\toprule
Decision policy & Threshold & Total cost (USD) & Cost / 1000 samples & Savings vs.\ treat-none (\%) & Savings vs.\ best baseline (\%) \\
\midrule
Treat none (no intervention) & --- & 30{,}480 & 6{,}285.83 & 0.0 & $-19.0$ \\
Treat all (expedite every sample) & --- & 72{,}735 & 15{,}000.00 & $-138.6$ & $-184.1$ \\
Priority rule (expedite STAT only) & --- & 25{,}605 & 5{,}280.47 & 16.0 & 0.0 \\
Model @ threshold 0.50 & 0.50 & 21{,}900 & 4{,}516.40 & 28.1 & 14.5 \\
Model @ cost-optimal threshold & 0.62 & 19{,}545 & 4{,}030.73 & \textbf{35.9} & \textbf{23.7} \\
\bottomrule
\end{tabular}
\end{table*}

\subsection{When the stakes rise: sensitivity}
The value of the framework grows with the cost of failure. As the ratio
$C_{\mathrm{B}}/C_{\mathrm{I}}$ increases from 2 to 50, the cost-optimal threshold
falls, the system intervenes more aggressively, and savings relative to doing
nothing rise from 6.1\% to 77.9\%. This is precisely the regime of time-critical
donor screening, where a delayed result can forfeit organ viability and the
effective $C_{\mathrm{B}}$ is very large. The same analytics that yield modest
savings in routine chemistry become decisive in high-stakes biomedical operations,
which is where we argue the approach has the greatest strategic value.

\section{Discussion}\label{sec:discussion}
The central message is that the managerial payoff lies not in prediction alone but
in the disciplined translation of prediction into a cost-rational, explainable
decision. Three implications follow for practice. First, the dominant,
controllable drivers of TAT breach are operational: workload, analyzer
utilization, shift, and transport, so the most effective interventions are
managerial (dynamic staffing, capacity planning, transport scheduling) rather than
purely technical. Second, the asymmetry between the two tasks is itself a decision:
because rejection resists prediction while TAT breach does not, scarce analytics
and automation capacity should be directed at TAT management and at upstream
process control for specimen quality. Third, the sensitivity analysis gives
managers a principled way to set the operating point according to the stakes of
their setting, rather than defaulting to an arbitrary threshold.

For diagnostic enterprises and organ-procurement organizations pursuing digital
transformation, the framework offers a template that is transparent enough to
govern and auditable enough to defend, properties that matter in regulated
environments where an unexplained model is difficult to adopt. It also formalizes,
in the open literature, a class of operational analytics that large testing
networks often develop privately, lowering the barrier for smaller laboratories to
adopt evidence-based decision support.

\section{Limitations}\label{sec:limits}
This study has clear limitations. The dataset is synthetic and literature-calibrated;
although this makes the work fully reproducible and avoids privacy constraints, the
absolute performance and cost figures should be read as illustrative of the
methodology rather than as estimates for any specific laboratory. The cost model
is deliberately simple: linear unit costs and an assumption that intervention
mitigates a true breach, and real settings involve partial mitigation, capacity
constraints on how many samples can be expedited, and non-linear failure costs.
Prospective validation on de-identified operational data under an appropriate
data-use agreement is the essential next step, together with extension of the cost
layer to capacitated and dynamic intervention policies.

\section{Conclusion}\label{sec:conclusion}
We presented an explainable decision-intelligence framework for clinical
laboratory operations that moves beyond turnaround-time prediction to
cost-rational, auditable decision-making. On a reproducible, literature-calibrated
benchmark, the approach predicts TAT breach with strong discrimination, attributes
risk to managerially-actionable drivers through SHAP, and converts predictions into
expedite decisions that cut expected operational cost by up to 35.9\% against doing
nothing and 23.7\% against the best managerial heuristic, with the advantage
widening sharply as the stakes rise toward time-critical donor screening. By
uniting prediction, explanation, and economic optimization, the framework offers
laboratory and diagnostic-service leaders a practical bridge from analytics to
governed decisions, and a replicable template for decision-intelligence research
in regulated biomedical operations.

\section*{Reproducibility}
The data generator, experiment pipeline, explainability and cost-analysis code,
and all figures and tables are available at
\url{https://github.com/jaisrithad/clinical-lab-decision-intelligence}.
All results are reproducible from a fixed random seed.

\section*{Declaration of Competing Interest}
The author declares that she has no known competing financial interests or
personal relationships that could have appeared to influence the work reported in
this paper.

\section*{Funding}
This research did not receive any specific grant from funding agencies in the
public, commercial, or not-for-profit sectors.

\section*{Data Availability}
All data used in this study are synthetic and generated by the released code; no
human-subjects, patient, or donor data were used. The data generator and the
generated dataset are available in the public repository linked above, enabling
full reproduction of every reported result.

\section*{Ethics Statement}
This study did not involve human participants, animals, or identifiable personal
data. No ethical approval was therefore required.

\section*{Declaration of Generative AI and AI-Assisted Technologies}
During the preparation of this work, the author used generative AI and
AI-assisted tools to support drafting, code development, and language editing.
After using these tools, the author reviewed and edited the content as needed and
takes full responsibility for the content of the publication.

\begin{thebibliography}{1}

\bibitem{tatml2022}
E.~R.~Tsai, D.~Demirtas, N.~Hoogendijk, A.~N.~Tintu, and R.~J.~Boucherie,
``Turnaround time prediction for clinical chemistry samples using machine
learning,'' \emph{Clinical Chemistry and Laboratory Medicine}, vol.~60, no.~12,
pp.~1902--1910, 2022.

\bibitem{pathoml2026}
W.~Kinnares, ``Predicting anatomic pathology turnaround time using machine
learning models,'' \emph{Journal of the Thai Medical Informatics Association},
vol.~12, no.~1, pp.~49--58, 2026.

\bibitem{orange2026}
Y.~Goren and D.~I.~Topcu, ``Application of data mining for monitoring turnaround
time in clinical laboratories: a workflow-based model using Orange visual
programming,'' \emph{International Journal of Medical Biochemistry}, vol.~9, no.~3,
pp.~162--169, 2026.

\bibitem{urinepretat2025}
S.~Lv, H.~Ye, Y.~Li, and J.~Zhang, ``Applying a logistic regression--clustering
joint model to analyze the causes of prolonged pre-analytic turnaround time for
urine culture testing in hospital wards,'' \emph{Frontiers in Digital Health},
vol.~7, 2025.

\bibitem{plebani2010}
M.~Plebani, ``The detection and prevention of errors in laboratory medicine,''
\emph{Annals of Clinical Biochemistry}, vol.~47, no.~2, pp.~101--110, 2010.

\bibitem{hawkins2012}
R.~Hawkins, ``Managing the pre- and post-analytical phases of the total testing
process,'' \emph{Annals of Laboratory Medicine}, vol.~32, no.~1, pp.~5--16, 2012.

\bibitem{clsi2021}
Clinical and Laboratory Standards Institute, \emph{Collection, Transport, and
Processing of Specimens for Clinical Testing}, CLSI guideline GP41, CLSI, 2021.

\bibitem{crosby1979}
P.~B.~Crosby, \emph{Quality is Free: The Art of Making Quality Certain}.
New York, NY, USA: McGraw-Hill, 1979.

\bibitem{shap2017}
S.~M.~Lundberg and S.-I.~Lee, ``A unified approach to interpreting model
predictions,'' in \emph{Advances in Neural Information Processing Systems (NeurIPS)},
2017, pp.~4765--4774.

\bibitem{breiman2001}
L.~Breiman, ``Random forests,'' \emph{Machine Learning}, vol.~45, no.~1,
pp.~5--32, 2001.

\bibitem{xgboost2016}
T.~Chen and C.~Guestrin, ``XGBoost: a scalable tree boosting system,'' in
\emph{Proc.\ 22nd ACM SIGKDD Int.\ Conf.\ Knowledge Discovery and Data Mining},
2016, pp.~785--794.

\bibitem{sklearn2011}
F.~Pedregosa \emph{et al.}, ``Scikit-learn: machine learning in Python,''
\emph{Journal of Machine Learning Research}, vol.~12, pp.~2825--2830, 2011.

\end{thebibliography}

\end{document}
